%% Chapitre 12 — Pièges, cas vécus et comment s'en sortir.
%% RÈGLE : uniquement des commandes sémantiques bookforge.
%% Aucun terme du glossaire n'est introduit ici (tous déjà vus aux ch. précédents).

\chapitre{Pièges, cas vécus et comment s'en sortir}

\introChapitre{
  Ce chapitre rassemble les histoires qui font mal : l'agent qui boucle, qui
  supprime trop, qui ment avec aplomb, qui part dans le décor. Pour chaque
  piège, un signal d'alerte, une cause, et une manœuvre de récupération. C'est
  l'expérience qu'on n'a pas envie d'acquérir au prix fort.
}

\paragraphe{Tu sais maintenant vérifier le travail de l'agent et poser des
garde-fous. Reste la partie que personne n'aime raconter : les sessions qui
dérapent. Bonne nouvelle, les dérapages se répètent. Une poignée de scénarios
revient sans cesse, et chacun a sa parade. Apprends à les reconnaître tôt,
quand le coût de la sortie est encore faible.}

\section{L'agent qui boucle ou s'enlise}

\paragraphe{Tu lui demandes de faire passer un test. Il modifie, lance, échoue,
re-modifie, relance, échoue encore. Trois tentatives, cinq, dix. Chaque essai
défait à moitié le précédent. La barre d'avancement tourne, le compteur de
jetons grimpe, et le test reste rouge.}

\begin{avertissement}
Signal d'alerte : l'agent répète le même geste sur le même fichier sans que le
résultat change. Si tu vois passer deux fois \emphase{presque} le même diff,
ou la même erreur reformulée, il ne progresse plus, il s'agite.
\end{avertissement}

\paragraphe{La cause est presque toujours la même : il lui manque une
information décisive, et il ne le sait pas. Le test échoue pour une raison qui
n'est pas dans son contexte (une variable d'environnement absente, une
dépendance non installée, un comportement attendu qu'il a mal deviné). Plutôt
que de signaler son ignorance, il tente des variations au hasard. Un agent ne
dit jamais spontanément « je suis bloqué » : il essaie encore.}

\begin{astuce}
Manœuvre : coupe, puis change de registre. Interromps la boucle et demande un
\emphase{diagnostic}, pas une correction : « Ne modifie rien. Explique-moi
pourquoi ce test échoue et ce qui te manque pour le savoir. » Tu rends visible
l'information absente, tu la lui fournis, et tu relances. Neuf fois sur dix, la
correction tient du premier coup.
\end{astuce}

\paragraphe{Fixe-toi une règle simple : \important{deux échecs identiques, tu
reprends la main.} Au troisième tour à l'aveugle, l'agent ne raisonne plus, il
brûle des jetons.}

\section{Suppressions et modifications excessives}

\paragraphe{Tu avais demandé de nettoyer un fichier. Tu relis le diff trop
vite, tu acceptes. Cinq minutes plus tard, une fonction utilisée ailleurs a
disparu, des commentaires importants sont partis avec le « bruit », et un bloc
entier a été « simplifié » jusqu'à changer de comportement.}

\begin{avertissement}
Signal d'alerte : un diff dont le volume de suppressions dépasse largement ce
que la tâche justifiait. « Nettoie cet import » qui efface quarante lignes
mérite un arrêt, pas un clic d'approbation.
\end{avertissement}

\paragraphe{La cause : l'agent optimise pour la consigne littérale sans
mesurer les effets de bord. « Nettoyer » est ambigu ; il l'a interprété
largement. Et un périmètre flou produit toujours des dégâts larges. C'est
exactement ce que les garde-fous du chapitre précédent servent à contenir.}

\begin{astuce}
Manœuvre : si tu travailles sur une branche (et tu devrais), la sortie est
indolore. Tu jettes les changements non commités et tu repars :
\end{astuce}

\begin{extraitcode}{bash}
# Annuler tout ce qui n'est pas encore commité, fichiers suivis
git restore .

# Voir l'ampleur des dégâts AVANT de décider quoi garder
git diff --stat

# Si un commit est déjà passé, revenir au précédent sans rien perdre
git reset --soft HEAD~1
\end{extraitcode}

\paragraphe{La vraie leçon est en amont : commite \important{avant} de lancer
une tâche large, pas après. Un point de retour propre transforme une
catastrophe en simple \code{git restore}.}

\section{Le code confiant mais faux}

\paragraphe{C'est le piège le plus sournois, parce qu'il ne ressemble pas à un
piège. L'agent te livre une fonction nette, bien commentée, qui appelle une
méthode de bibliothèque au nom parfaitement crédible. Le ton est assuré.
Le problème : cette méthode n'existe pas. Ou elle existe, mais ne fait pas ce
qu'il affirme.}

\begin{avertissement}
Signal d'alerte : une explication \emphase{trop} fluide, sans la moindre
réserve, sur du code que tu n'as pas exécuté. L'aplomb n'est pas une preuve.
Une hallucination se présente toujours avec autant d'assurance qu'une vérité.
\end{avertissement}

\paragraphe{La cause tient à la nature du modèle : il produit ce qui est
\emphase{plausible}, pas ce qui est \emphase{vérifié}. Une signature
d'\code{API} inventée est statistiquement vraisemblable, donc elle sort avec la
même confiance que le reste. L'agent ne ment pas au sens humain ; il comble un
vide avec ce qui ressemble à la bonne réponse.}

\begin{astuce}
Manœuvre : ne discute pas, exécute. Le seul antidote fiable, c'est de faire
tourner le code (ou les tests). Demande à l'agent de lancer lui-même la
vérification : « Écris un test minimal qui prouve que cette fonction marche, et
exécute-le. » Tu remplaces une affirmation par une preuve. Si la méthode est
fantôme, l'erreur d'exécution la démasque immédiatement.
\end{astuce}

\paragraphe{Méfie-toi en particulier des API récentes, des bibliothèques de
niche et des versions précises : c'est là que le modèle invente le plus, faute
d'avoir assez vu ces cas.}

\section{Dérive de contexte}

\paragraphe{La session a commencé sur une migration de base de données. Deux
heures et trente échanges plus tard, tu lui demandes un détail, et il te répond
à côté : il a oublié une décision prise au début, il réintroduit un bug que
vous aviez corrigé ensemble, il confond deux fichiers. L'agent du début n'est
plus tout à fait là.}

\begin{avertissement}
Signal d'alerte : l'agent contredit une consigne donnée plus tôt dans la même
session, ou « redécouvre » un problème déjà résolu. Sa mémoire de la
conversation s'effrite par le début.
\end{avertissement}

\paragraphe{La cause : la fenêtre de contexte sature. Quand elle déborde, les
échanges les plus anciens sont écartés ou compressés, et la qualité se dégrade
en silence. Une session trop longue, trop touffue, finit toujours par perdre le
fil de ses propres décisions.}

\begin{astuce}
Manœuvre : redémarre proprement plutôt que d'insister. Ouvre une nouvelle
session, et rouvre-la avec un résumé court de l'état : ce qui est fait, ce qui
reste, les décisions à ne pas rejouer. Mieux encore, consigne les décisions
durables dans la mémoire projet pour qu'elles survivent à n'importe quelle
session.
\end{astuce}

\paragraphe{Préviens la dérive : une session, une mission. Quand la tâche
change vraiment de nature, repars à neuf plutôt que d'empiler.}

\section{Sur-confiance du développeur}

\paragraphe{Le piège le plus dangereux n'est pas dans l'agent. Il est dans toi.
Après dix bons diffs d'affilée, tu cesses de lire. Tu cliques « accepter » au
réflexe. Tu passes en pilote automatique. Et c'est précisément le onzième diff,
celui qui touche l'authentification, qui contenait l'erreur que tu n'as pas
vue.}

\begin{avertissement}
Signal d'alerte : tu approuves plus vite que tu ne lis. Si tu ne pourrais pas
résumer ce que le dernier diff a changé, tu ne l'as pas relu, tu l'as gobé.
\end{avertissement}

\paragraphe{La cause est humaine : une série de réussites endort la vigilance.
Plus l'agent est bon, plus le risque monte, parce que la confiance se
construit là où elle coûtera le plus cher. La vérification n'est pas un manque
de confiance envers l'agent ; c'est ta part du travail, et elle ne se délègue
pas.}

\begin{astuce}
Manœuvre : module ton attention selon l'enjeu. Un changement de logique métier,
de sécurité ou de migration de données se relit ligne par ligne, toujours, même
au cinquantième diff. Le confort de routine se réserve aux zones où une erreur
se rattrape sans douleur.
\end{astuce}

\section{Boîte à outils de récupération}

\paragraphe{Quand une session dérape, tu n'as pas besoin de paniquer, tu as
besoin de réflexes. Voici les gestes qui sortent de presque toutes les
situations vues dans ce chapitre :}

\begin{liste}
  \element{\important{Couper.} Interromps l'agent dès que tu vois une boucle ou
  un dérapage. Chaque tour supplémentaire ne fait qu'aggraver et coûter.}
  \element{\important{Annuler.} \code{git restore} pour les changements non
  commités, \code{git reset} pour revenir à un commit propre. Travailler sur
  une branche rend tout ça réversible.}
  \element{\important{Recadrer.} Passe de la correction au diagnostic : « ne
  modifie rien, explique-moi pourquoi ». Tu rends visible l'information qui
  manquait.}
  \element{\important{Prouver.} Pour tout code dont tu doutes, exige un test
  exécuté, pas une explication. Une preuve vaut mille assurances.}
  \element{\important{Repartir.} Quand le contexte est saturé ou pollué, une
  session neuve avec un résumé court bat dix relances sur une session fatiguée.}
\end{liste}

\resumeChapitre{
  \element{\important{Agent qui boucle} : deux échecs identiques, tu reprends la main — interromps et demande un diagnostic plutôt qu'une correction, pour rendre visible l'information qui manque.}
  \element{\important{Suppressions excessives} : un diff dont les suppressions dépassent la tâche s'arrête avant approbation ; commite avant de lancer pour que \code{git restore} soit toujours disponible.}
  \element{\important{Code confiant mais faux} : ne discute pas, exécute — fais tourner un test minimal plutôt que d'accepter une explication fluide.}
  \element{\important{Dérive de contexte} : quand l'agent contredit une décision prise plus tôt, redémarre avec un résumé court plutôt que d'insister sur une session saturée.}
  \element{\important{Sur-confiance du développeur} : module ton attention selon l'enjeu — sécurité, logique métier et migrations se relisent ligne par ligne, même au cinquantième diff.}
}

\paragraphe{Aucun de ces pièges ne disparaîtra : ils sont inhérents à la
nature de l'agent. Mais chacun se désamorce avec le bon réflexe au bon moment.
La différence entre une session pénible et une session productive ne tient pas
à la chance ; elle tient à ces habitudes. \important{C'est justement
d'habitudes qu'on va parler au dernier chapitre} : comment tisser tout ce que
tu as appris dans ta journée de travail, pour que ces réflexes deviennent
naturels plutôt que durement appris.}
